\scanpage{10}
160.

\uline{Proof} (i) $\Leftrightarrow$ (i)'. $R$ can be desingularized if its completion $\widehat R$ is normal (Remark 16.2); hence the equivalence of (i) and (i)' is given by 17.3 and 17.2.

(i) $\Rightarrow$ (ii). By 17.3 and 17.2, (i) implies that $R$ has a rational singularity and that $H(R)$ is trivial. Since $R/m$ is algebraically closed, ``1'' is the only integer which can appear in the configuration diagram on a minimal desingularization of $R$. Proposition 22.5 shows then that $R$ has multiplicity two. We must therefore be in CASE Vd of \S24, so that for a suitable basis $\{x,y,z\}$ of $m$ there is a relation
\[
z^2+\alpha y^3+tx^4y+\beta x^5+px^3z+qx^2y^2+rxyz+sy^2z=0
\qquad\ldots(6)
\]
where $\alpha$ and $\beta$ are units in $R$, and $p,q,r,s,t$ are in $R$ (this is derived from equation (5)' in \S24, where $\rho$ is a \uline{non-unit} i.e. $\rho\in(x,y,z)R$; also we have put $\beta$ in place of $\sigma$).

Setting
\[
z'=z+(p/2)x^3+(r/2)xy+(s/2)y^2
\]
we have
\[
(z')^2+\alpha'y^3+t'x^4y+\beta'x^5+q'x^2y^2=0
\]
for suitable $\alpha',t',\beta',q'$, with $\alpha'\equiv\alpha$, $\beta'\equiv\beta$ (mod. $m$). In other words, for suitable $x,y,z$ we may assume in (6) that $p=r=s=0$.
\mspage{9}
